\section{法的演进}
\begin{enumerate}
  \item 不重要。
  \item 法律产生的标志。
  \item 法律、国家没有先后之分，同时产生，同步发展。
  \item 法律意识。对于法的观点、知识和心理态度的总称。是一国法律传统中相对比较稳定的部分。
  \item 法的历史类型。经济基础、阶级本质。
  \item 法的继承。\zd{不同历史类型}的法律制度之间的延续和继受。表现为\zd{新旧法律}承接和继受。
  \item 法的移植。\zd{同时代}，国外制度的吸收和学习。包括外国法律、国际法律和惯例。
  \item 法律传统。family、阵营。从意识形态分。
        \begin{enumerate}
          \item 资本主义法系。沿袭哪个\zd{历史传统}。
                \begin{enumerate}
                  \item 民法法系/大陆法系/罗马法系/法德法系/法典法系/职权主义纠问制。
                  \item 普通法系/海洋法系/英美法系/判例法系/当事人对抗制。
                \end{enumerate}
          \item 社会主义法系。
        \end{enumerate}
  \item 法的现代化。物质到精神、制度到观念的\zd{整体社会变迁}。
        \begin{enumerate}
          \item 类型。动力来自哪。内发现代化 \yjt 自下而上、外源现代化 \yjt 自上而下。
          \item 清末修律开始，中国法的现代化一直是\zd{立法主导}性。自上而下建立。
        \end{enumerate}
\end{enumerate}
